%==================================================================
% FILE UTAMA — PROPOSAL TUGAS AKHIR
% Politeknik Negeri Banyuwangi
% Acuan: Pedoman Mutu Tugas Akhir Poliwangi 2023, Tabel 6.3
%
% Berkas ini dipisah dari Laporan supaya mudah dipakai sendiri-sendiri.
% Folder proposal/ sudah lengkap: identitas, bab, gambar, kode, referensi.
%
% Cara memakai:
%   1. Isi identitas di  a0-identitas.tex
%   2. Tulis isi bab di  b1-bab1.tex sampai b3-bab3.tex
%   3. Tambah referensi di  a9-pustaka.bib
%   4. Kompilasi: pdflatex -> bibtex -> pdflatex -> pdflatex
%==================================================================

\documentclass[12pt,a4paper,oneside,fleqn]{report}

%%========================================================================
%% IDENTITAS TUGAS AKHIR
%% Isi semua bagian di bawah ini. Bagian yang masih contoh WAJIB diganti.
%%========================================================================

%--- Judul -----------------------------------------------------------------
% Ditulis singkat (sekitar 20 kata), kapital, membentuk piramida terbalik.
% Di LaTeX cukup tulis biasa; pemenggalan baris diatur otomatis.
\providecommand{\judulid}{Rancang Bangun Sistem Deteksi Kesegaran Ikan Berbasis Pengolahan Citra Digital dan Jaringan Saraf Tiruan}
\providecommand{\judulen}{Design of a Fish Freshness Detection System Based on Digital Image Processing and Artificial Neural Networks}

%--- Mahasiswa -------------------------------------------------------------
\providecommand{\penulis}{Nama Lengkap Mahasiswa}
\providecommand{\nim}{3210xxxxxx}

%--- Program studi ---------------------------------------------------------
\providecommand{\jenjang}{Sarjana Terapan}          % atau: Diploma Tiga
\providecommand{\gelar}{Sarjana Terapan (S.Tr.)}    % atau: Ahli Madya (A.Md)
\providecommand{\prodi}{Teknologi Rekayasa Perangkat Lunak}
\providecommand{\jurusan}{Teknik Informatika}
\providecommand{\institusi}{Politeknik Negeri Banyuwangi}
\providecommand{\kota}{Banyuwangi}
\providecommand{\tahun}{\the\year}

%--- Pembimbing dan penguji ------------------------------------------------
\providecommand{\pembimbingA}{Nama Pembimbing I, S.Kom., M.Kom.}
\providecommand{\NIDNpembimbingA}{00xxxxxxxx}
\providecommand{\pembimbingB}{Nama Pembimbing II, S.T., M.T.}
\providecommand{\NIDNpembimbingB}{00xxxxxxxx}

\providecommand{\pengujiA}{Nama Penguji I, S.Kom., M.Kom.}
\providecommand{\pengujiB}{Nama Penguji II, S.T., M.T.}

%--- Pejabat penanda tangan ------------------------------------------------
\providecommand{\ketuajurusan}{Nama Ketua Jurusan, S.T., M.T.}
\providecommand{\NIPketuajurusan}{19xxxxxxxxxxxxxxxx}
\providecommand{\kaprodi}{Nama Ketua Program Studi, S.Kom., M.Kom.}
\providecommand{\NIPkaprodi}{19xxxxxxxxxxxxxxxx}

%--- Pelaksanaan seminar / sidang ------------------------------------------
\providecommand{\hariujian}{Senin}
\providecommand{\tanggalUjian}{1 Januari \tahun}
\providecommand{\tempatujian}{Ruang Sidang Jurusan \jurusan}

%--- Warna sampul sesuai program studi (Pedoman Tabel 6.2) -----------------
% Teknik Informatika   : kuning, tulisan hitam   -> FFFF00
% Teknik Sipil         : biru tua, tulisan emas  -> 000080
% Teknik Mesin         : merah, tulisan emas     -> FF0000
% Manajemen Bisnis Pariwisata : oranye           -> FF7F00
% Agribisnis           : hijau tua               -> 009F3C
% Teknologi Pengolahan Hasil Ternak : coklat tua -> 964B00
% Teknik Manufaktur Kapal : biru muda            -> 0080FF
\providecommand{\warnasampul}{FFFF00}

%==================================================================
% PENGATURAN UTAMA DOKUMEN
% Template Tugas Akhir Politeknik Negeri Banyuwangi
% Acuan: Pedoman Mutu Tugas Akhir Poliwangi 2023, Bab 6
%
% FILE INI TIDAK PERLU DIUBAH.
% Kalau ingin mengubah format, ubah dengan sadar dan catat alasannya.
%==================================================================

%------------------------------------------------------------------
% 1. Bahasa dan pengkodean huruf
%------------------------------------------------------------------
\usepackage[T1]{fontenc}
\usepackage[utf8]{inputenc}
\usepackage[indonesian]{babel}

%------------------------------------------------------------------
% 2. Huruf Times New Roman, jarak baris 1,5
%    (Pedoman 6.1 butir 2)
%------------------------------------------------------------------
\usepackage{mathptmx}     % huruf Times + rumus matematika bergaya Times
\usepackage{setspace}
\setstretch{1.5}          % jarak baris isi = 1,5 spasi

%------------------------------------------------------------------
% 3. Kertas A4 dan margin
%    (Pedoman Lampiran 1: kiri 3 cm, atas 2 cm, kanan 2 cm, bawah 2 cm)
%------------------------------------------------------------------
\usepackage[a4paper,
            left=3cm,
            top=2cm,
            right=2cm,
            bottom=2cm,
            headsep=0.6cm,
            footskip=1.0cm]{geometry}

%------------------------------------------------------------------
% 4. Paragraf: baris pertama masuk 1,25 cm, jarak sebelum/sesudah 0
%    (Pedoman 6.1 butir 3)
%------------------------------------------------------------------
\usepackage{indentfirst}
\setlength{\parindent}{1.25cm}
\setlength{\parskip}{0pt}

% Mengurangi masalah perataan baris pada teks terjustifikasi
\tolerance=200
\emergencystretch=2.5em
\hyphenpenalty=200
\hbadness=3000

%------------------------------------------------------------------
% 5. Judul BAB, sub-bab, dan sub-sub-bab
%    (Pedoman Lampiran 2: "BAB 1" lalu judul, rata tengah, tebal)
%------------------------------------------------------------------
\usepackage{titlesec}

\titleformat{\chapter}[display]
  {\normalfont\fontsize{12}{14}\bfseries\centering}
  {BAB \thechapter}
  {0.6\baselineskip}
  {\MakeUppercase}
\titlespacing*{\chapter}{0pt}{-1.4cm}{1.2\baselineskip}

\titleformat{\section}
  {\normalfont\fontsize{12}{14}\bfseries}
  {\thesection}{0.6em}{}
\titlespacing*{\section}{0pt}{12pt}{6pt}

\titleformat{\subsection}
  {\normalfont\fontsize{12}{14}\bfseries}
  {\thesubsection}{0.6em}{}
\titlespacing*{\subsection}{0pt}{10pt}{4pt}

% Judul bab yang tidak bernomor (DAFTAR ISI, ABSTRAK, dst.)
\titleformat{name=\chapter,numberless}[block]
  {\normalfont\fontsize{12}{14}\bfseries\centering}
  {}{0pt}
  {\MakeUppercase}
\titlespacing*{name=\chapter,numberless}{0pt}{-1.4cm}{1.2\baselineskip}

%------------------------------------------------------------------
% 6. Keterangan tabel dan gambar
%    Tabel : nomor + judul di ATAS, rata kiri, tebal, spasi 1
%    Gambar: nomor + judul di BAWAH, rata tengah, tebal, spasi 1
%    (Pedoman 6.1 butir 4 dan 6)
%------------------------------------------------------------------
\usepackage[font={normalsize,stretch=1},
            labelfont=bf,
            textfont=bf,
            labelsep=space,
            singlelinecheck=false]{caption}
\usepackage{subcaption}

\captionsetup[table]{position=above, justification=raggedright, skip=6pt}
\captionsetup[figure]{position=below, justification=centering, skip=6pt}
\captionsetup[lstlisting]{position=above, justification=raggedright, skip=6pt}

% Penomoran tabel/gambar mengikuti nomor bab: Tabel 2.1, Gambar 3.2
\renewcommand{\thetable}{\arabic{chapter}.\arabic{table}}
\renewcommand{\thefigure}{\arabic{chapter}.\arabic{figure}}

%------------------------------------------------------------------
% 7. Rumus matematika
%    Nomor rumus mengikuti nomor bab dan rata kanan (Pedoman 6.1 butir 5)
%    Rumus dimulai dari batas kiri margin (opsi fleqn di laporan.tex)
%------------------------------------------------------------------
\usepackage{amsmath}
\numberwithin{equation}{chapter}
\setlength{\mathindent}{0pt}

%------------------------------------------------------------------
% 8. Daftar isi
%------------------------------------------------------------------
\usepackage{tocloft}

\renewcommand{\cftchapfont}{\normalfont\bfseries}
\renewcommand{\cftchappagefont}{\normalfont\bfseries}
\renewcommand{\cftchappresnum}{BAB }
\renewcommand{\cftchapaftersnum}{ }
\newlength{\lebarbab}
\settowidth{\lebarbab}{\bfseries BAB 1\ }
\cftsetindents{chap}{0pt}{\lebarbab}

\setlength{\cftbeforechapskip}{0pt}
\setlength{\cftbeforesecskip}{0pt}
\setlength{\cftbeforesubsecskip}{0pt}
\renewcommand{\cftsecleader}{\cftdotfill{\cftdotsep}}
\setlength{\cftsecindent}{0pt}
\setlength{\cftsecnumwidth}{2.6em}
\setlength{\cftsubsecindent}{2.6em}
\setlength{\cftsubsecnumwidth}{3.2em}
\setlength{\cftfigindent}{0pt}
\setlength{\cfttabindent}{0pt}
\renewcommand{\cftdotsep}{1}

% Daftar isi, daftar tabel, dan daftar gambar memakai spasi 1
\renewcommand{\cfttoctitlefont}{\hfill\normalfont\bfseries}
\renewcommand{\cftaftertoctitle}{\hfill}
\renewcommand{\cftloftitlefont}{\hfill\normalfont\bfseries}
\renewcommand{\cftafterloftitle}{\hfill}
\renewcommand{\cftlottitlefont}{\hfill\normalfont\bfseries}
\renewcommand{\cftafterlottitle}{\hfill}

%------------------------------------------------------------------
% 9. Tabel dan gambar
%------------------------------------------------------------------
\usepackage{graphicx}
\graphicspath{{gambar/}}
\usepackage{float}
\usepackage{tabularx}
\usepackage{multirow}
\usepackage{longtable}
\usepackage{rotating}       % halaman landscape
\usepackage{array}

%------------------------------------------------------------------
% 10. Daftar butir
%------------------------------------------------------------------
\usepackage{enumitem}
\setlist{nosep}

\newenvironment{packed_enum}
  {\begin{enumerate}[leftmargin=1.5\parindent]}
  {\end{enumerate}}
\newenvironment{packed_item}
  {\begin{itemize}[leftmargin=1.5\parindent]}
  {\end{itemize}}

%------------------------------------------------------------------
% 11. Daftar pustaka gaya APA (penulis, tahun)
%------------------------------------------------------------------
\usepackage[round,authoryear]{natbib}
\bibliographystyle{apalike}

%------------------------------------------------------------------
% 12. Kode program
%------------------------------------------------------------------
\usepackage{listings}
\usepackage[table]{xcolor}   % opsi table: untuk \rowcolor dan \cellcolor

\definecolor{codegreen}{rgb}{0,0.5,0}
\definecolor{codegray}{rgb}{0.45,0.45,0.45}
\definecolor{codepurple}{rgb}{0.55,0,0.65}
\definecolor{backcolour}{rgb}{0.97,0.97,0.95}

\lstdefinestyle{poliwangi}{
    backgroundcolor=\color{backcolour},
    commentstyle=\color{codegreen},
    keywordstyle=\color{magenta}\bfseries,
    numberstyle=\tiny\color{codegray},
    stringstyle=\color{codepurple},
    basicstyle=\ttfamily\fontsize{9}{10.8}\selectfont,
    columns=fullflexible,
    breakatwhitespace=false,
    breaklines=true,
    captionpos=t,
    keepspaces=true,
    numbers=left,
    numbersep=6pt,
    showspaces=false,
    showstringspaces=false,
    showtabs=false,
    tabsize=2,
    frame=single,
    rulecolor=\color{codegray},
    xleftmargin=1.2em,
    framexleftmargin=1.2em
}
\lstset{style=poliwangi}

\AtBeginDocument{%
  \renewcommand{\lstlistingname}{Kode}%
  \renewcommand{\thelstlisting}{\arabic{chapter}.\arabic{lstlisting}}%
}

%------------------------------------------------------------------
% 13. Gambar dengan TikZ (diagram alir, grafik)
%------------------------------------------------------------------
\usepackage{tikz}
\usetikzlibrary{shapes.geometric, arrows.meta, positioning, calc}
\usepackage{pgfplots}
\pgfplotsset{compat=1.18}

\tikzset{
  kotak/.style   = {rectangle, rounded corners, minimum width=3cm,
                    minimum height=0.9cm, text centered, draw=black,
                    align=center, font=\small},
  proses/.style  = {kotak, fill=blue!8},
  mulai/.style   = {kotak, fill=green!12},
  keputusan/.style = {diamond, aspect=2.2, draw=black, align=center,
                      font=\small, fill=orange!12, inner sep=1pt},
  panah/.style   = {thick, -{Stealth[length=2.5mm]}}
}

%------------------------------------------------------------------
% 14. Tautan dan rujukan otomatis (Gambar 2.1, Tabel 3.2, persamaan 4.1)
%------------------------------------------------------------------
\usepackage[hidelinks]{hyperref}
\usepackage{cleveref}

\crefname{figure}{Gambar}{Gambar}
\Crefname{figure}{Gambar}{Gambar}
\crefname{table}{Tabel}{Tabel}
\Crefname{table}{Tabel}{Tabel}
\crefname{equation}{persamaan}{persamaan}
\Crefname{equation}{Persamaan}{Persamaan}
\crefname{lstlisting}{Kode}{Kode}
\Crefname{lstlisting}{Kode}{Kode}

%------------------------------------------------------------------
% 15. Perintah bantu
%------------------------------------------------------------------
\usepackage{ifthen}

% Halaman kosong diberi keterangan, bukan halaman benar-benar kosong
\newcommand{\halamankosong}{%
  \clearpage
  \thispagestyle{empty}%
  \vspace*{\fill}%
  \begin{center}\textit{``Halaman ini sengaja dikosongkan''}\end{center}%
  \vspace*{\fill}%
  \clearpage
}

% Jarak bantu: \jarak{2} = 2 x 1 spasi
\newcommand{\jarak}[1]{\vspace{#1\baselineskip}}

% Bagian awal (ABSTRAK, KATA PENGANTAR, dst.) dimasukkan ke daftar isi
% TANPA awalan "BAB".
\makeatletter
\newcommand{\l@front}[2]{%
  \par\noindent\cftchapfont #1\nobreak
  \leaders\hbox{$\m@th\mkern\cftdotsep mu\hbox{.}\mkern\cftdotsep mu$}%
  \hfill\nobreak{\cftchappagefont #2}\par}
\makeatother

\newcommand{\frontmatterentry}[1]{%
  \phantomsection
  \addcontentsline{toc}{front}{#1}%
}

% Lampiran: bernomor urut sendiri dan otomatis masuk DAFTAR LAMPIRAN
\newcounter{lampiran}
\makeatletter
\newcommand{\l@lampiran}[2]{%
  \par\noindent LAMPIRAN \hb@xt@0.6em{}#1\nobreak
  \leaders\hbox{$\m@th\mkern\cftdotsep mu\hbox{.}\mkern\cftdotsep mu$}%
  \hfill\nobreak #2\par}
\makeatother

\newcommand{\lampiran}[1]{%
  \clearpage
  \refstepcounter{lampiran}%
  \addcontentsline{loa}{lampiran}{\thelampiran\hspace{0.8em}#1}%
  \phantomsection
  \begin{center}
    \textbf{LAMPIRAN \thelampiran}\par
    \vspace{0.2\baselineskip}
    \textbf{#1}\par
  \end{center}
  \vspace{0.4\baselineskip}%
}

\makeatletter
\newcommand{\daftarlampiran}{%
  \chapter*{DAFTAR LAMPIRAN}%
  \frontmatterentry{DAFTAR LAMPIRAN}%
  \begin{spacing}{1}%
  \@starttoc{loa}%
  \end{spacing}%
}
\makeatother

%==================================================================
% AKHIR PENGATURAN
%==================================================================


% Dokumen ini adalah Proposal
\providecommand{\Jenis}{Proposal}

\ifthenelse{\equal{\Jenis}{Proposal}}{%
  \newcommand{\juduljenis}{PROPOSAL TUGAS AKHIR}%
  \newcommand{\jenisjudul}{Proposal Tugas Akhir}%
}{%
  \newcommand{\juduljenis}{LAPORAN TUGAS AKHIR}%
  \newcommand{\jenisjudul}{Laporan Tugas Akhir}%
}

\begin{document}

%==================================================================
% BAGIAN AWAL PROPOSAL
%==================================================================
%==================================================================
% SAMPUL MUKA / SOFTCOVER
% Acuan: Pedoman Lampiran 4
%
% Ingin sampul berwarna sesuai program studi? Hapus tanda % pada
% empat baris tikz di bawah (butuh dua kali kompilasi).
%==================================================================
\begin{titlepage}
\thispagestyle{empty}%
\centering

%--- Latar berwarna (opsional) -------------------------------------
% \definecolor{warnasampul}{HTML}{\warnasampul}
% \begin{tikzpicture}[remember picture, overlay]
%   \fill[warnasampul] (current page.south west) rectangle (current page.north east);
% \end{tikzpicture}%
%-------------------------------------------------------------------

\vspace*{0.5cm}

%--- Judul: 14 pt, tebal, kapital, spasi 1 -------------------------
\begin{spacing}{1}
\begin{center}
{\fontsize{14}{16.8}\selectfont\bfseries \MakeUppercase{\judulid}\par}
\end{center}
\end{spacing}

\vspace{4\baselineskip}

%--- Jenis dokumen: 12 pt, tebal, spasi 1 --------------------------
\begin{spacing}{1}
\begin{center}
{\fontsize{12}{14.4}\selectfont\bfseries \juduljenis\par}
\end{center}
\end{spacing}

\vspace{2\baselineskip}

%--- Lambang Poliwangi 4 cm x 4 cm --------------------------------
\includegraphics[width=4cm,height=4cm,keepaspectratio]{logo-poliwangi.png}

\vspace{3\baselineskip}

%--- Nama dan NIM: 12 pt, tebal, spasi 1 --------------------------
\begin{spacing}{1}
\begin{center}
{\fontsize{12}{14.4}\selectfont\bfseries
 Oleh:\par
 \penulis\par
 NIM. \nim\par}
\end{center}
\end{spacing}

\vspace{2\baselineskip}

%--- Program studi dan tahun: 16 pt, tebal, spasi 1 ----------------
\begin{spacing}{1}
\begin{center}
{\fontsize{16}{19.2}\selectfont\bfseries
 PROGRAM STUDI \MakeUppercase{\prodi}\par
 \MakeUppercase{\institusi}\par
 \tahun\par}
\end{center}
\end{spacing}

\end{titlepage}

%==================================================================
% HALAMAN JUDUL (sampul dalam)
% Acuan: Pedoman Lampiran 5
%==================================================================
\begin{titlepage}
\thispagestyle{empty}%
\centering

\vspace*{0.3cm}

%--- Judul: 14 pt, tebal, kapital, spasi 1,5 -----------------------
\begin{spacing}{1}
\begin{center}
{\fontsize{14}{16.8}\selectfont\bfseries \MakeUppercase{\judulid}\par}
\end{center}
\end{spacing}

\vspace{2\baselineskip}

%--- Jenis dokumen: 12 pt, tebal, spasi 1,5 -----------------------
\begin{spacing}{1}
\begin{center}
{\fontsize{12}{14.4}\selectfont\bfseries \juduljenis\par}
\end{center}
\end{spacing}

\vspace{2\baselineskip}

\includegraphics[width=4cm,height=4cm,keepaspectratio]{logo-poliwangi.png}

\vspace{2\baselineskip}

%--- Kalimat pengantar, spasi 1,5 ---------------------------------
\begin{spacing}{1}
\begin{center}
{\fontsize{12}{14.4}\selectfont
 \ifthenelse{\equal{\Jenis}{Proposal}}{%
   Proposal Tugas Akhir ini dibuat dan diajukan sebagai salah satu
   syarat untuk melaksanakan Tugas Akhir pada Program Studi
   \jenjang\ \prodi\ \institusi.\par
 }{%
   Tugas Akhir ini dibuat dan diajukan untuk memenuhi salah satu
   syarat kelulusan Program Studi \jenjang\ \prodi\ dan mencapai
   gelar \gelar.\par
 }}
\end{center}
\end{spacing}

\vspace{2\baselineskip}

%--- Nama dan NIM: 12 pt, tebal, spasi 1,5 ------------------------
\begin{spacing}{1}
\begin{center}
{\fontsize{12}{14.4}\selectfont\bfseries
 Oleh:\par
 \penulis\par
 NIM. \nim\par}
\end{center}
\end{spacing}

\vspace{2\baselineskip}

%--- Program studi dan tahun: 16 pt, tebal ------------------------
\begin{spacing}{1}
\begin{center}
{\fontsize{16}{19.2}\selectfont\bfseries
 PROGRAM STUDI \MakeUppercase{\prodi}\par
 \MakeUppercase{\institusi}\par
 \tahun\par}
\end{center}
\end{spacing}

\end{titlepage}


% Halaman judul dihitung sebagai halaman i tetapi tidak dicetak,
% sehingga Lembar Pengesahan bernomor ii
\pagenumbering{roman}
\setcounter{page}{2}

%==================================================================
% LEMBAR PENGESAHAN PROPOSAL TUGAS AKHIR
% Acuan: Pedoman Lampiran 6
%==================================================================
\chapter*{LEMBAR PENGESAHAN}
\thispagestyle{plain}
\frontmatterentry{LEMBAR PENGESAHAN}
\begin{center}\textbf{PROPOSAL TUGAS AKHIR}\end{center}

\vspace{0.5\baselineskip}

\noindent
\begin{tabular}{@{}l@{\hspace{0.6em}}c@{\hspace{0.6em}}p{11.4cm}@{}}
Judul  & : & \judulid \\[2pt]
Oleh   & : & \penulis \\[2pt]
NIM.   & : & \nim \\
\end{tabular}

\vspace{2.5\baselineskip}

\begin{center}
\textbf{Telah diuji pada :}
\end{center}

\vspace{0.5\baselineskip}

\noindent
\begin{tabular}{@{}l@{\hspace{0.6em}}c@{\hspace{0.6em}}p{11.4cm}@{}}
Hari    & : & \hariujian \\[2pt]
Tanggal & : & \tanggalUjian \\[2pt]
Tempat  & : & \tempatujian \\
\end{tabular}

\vspace{2.5\baselineskip}

\begin{center}
\textbf{Mengetahui / Menyetujui :}
\end{center}

\vspace{1\baselineskip}

\noindent
\begin{tabular}{@{}p{7.6cm}p{7.6cm}@{}}
Dosen Penguji : & Dosen Pembimbing : \\[3.5cm]
1.\ \makebox[4.4cm]{\dotfill} & 1.\ \makebox[4.4cm]{\dotfill} \\
\phantom{1.}\ NIK/NIP. & \phantom{1.}\ NIK/NIP. \\[2.2cm]
2.\ \makebox[4.4cm]{\dotfill} & 2.\ \makebox[4.4cm]{\dotfill} \\
\phantom{1.}\ NIK/NIP. & \phantom{1.}\ NIK/NIP. \\
\end{tabular}

\clearpage

%==================================================================
% ABSTRAK (Bahasa Indonesia)
% Acuan: Pedoman Lampiran 9A
%   - judul TA 14 pt tebal, rata tengah, spasi 1
%   - nama, NIM, pembimbing 2 spasi di bawah judul
%   - kata ABSTRAK 2 spasi di bawah nama, 14 pt tebal
%   - isi abstrak spasi 1, baris pertama masuk 1,25 cm
%   - maksimal 350 kata, kata kunci 3-5 kata
%==================================================================
\clearpage
\thispagestyle{plain}
\frontmatterentry{ABSTRAK}

\begin{spacing}{1}
\begin{center}
{\fontsize{14}{16.8}\selectfont\bfseries \MakeUppercase{\judulid}\par}
\end{center}

\vspace{2\baselineskip}

\noindent
\begin{tabular}{@{}l@{\hspace{0.6em}}c@{\hspace{0.6em}}p{10.8cm}@{}}
Nama mahasiswa & : & \penulis \\
NIM            & : & \nim \\
Pembimbing     & : & 1.\ \pembimbingA \\
               &   & 2.\ \pembimbingB \\
\end{tabular}

\vspace{2\baselineskip}

\begin{center}
{\fontsize{14}{16.8}\selectfont\bfseries ABSTRAK\par}
\end{center}

\vspace{2\baselineskip}

%--- Isi abstrak: ganti dengan ringkasan penelitian Anda -------------
Penurunan mutu kesegaran ikan terjadi sangat cepat setelah ikan
ditangkap, sehingga penilaian kesegaran secara cepat dan objektif
menjadi kebutuhan penting bagi pembeli maupun pelaku usaha perikanan.
Selama ini penilaian kesegaran masih dilakukan secara organoleptik
oleh tenaga terlatih sehingga hasilnya bersifat subjektif dan sulit
dilakukan dalam jumlah besar. Penelitian ini bertujuan merancang
bangun sistem deteksi kesegaran ikan berbasis pengolahan citra digital
dan jaringan saraf tiruan yang mampu mengklasifikasikan tingkat
kesegaran ikan secara otomatis. Citra mata dan insang ikan diambil
menggunakan kamera digital pada kotak pengambilan gambar dengan
pencahayaan tetap, kemudian diolah melalui tahap penyekalaan ukuran,
perbaikan kontras, dan ekstraksi ciri warna pada ruang warna HSV.
Ciri yang diperoleh digunakan sebagai masukan jaringan saraf tiruan
dengan arsitektur perambatan balik dua lapisan tersembunyi. Pengujian
dilakukan terhadap 300 citra ikan yang terdiri atas tiga kelas
kesegaran, yaitu segar, cukup segar, dan tidak segar. Hasil pengujian
menunjukkan akurasi klasifikasi sebesar 93,33\% dengan waktu
pemrosesan rata-rata 0,42 detik per citra. Nilai presisi dan perolehan
tertinggi diperoleh pada kelas ikan segar, sedangkan kesalahan
terbanyak terjadi antara kelas cukup segar dan tidak segar karena
kemiripan warna insang. Dengan demikian sistem yang dirancang dapat
digunakan sebagai alternatif penilaian kesegaran ikan secara cepat,
objektif, dan tidak merusak.

\vspace{3\baselineskip}

\noindent
\textbf{Kata kunci:} jaringan saraf tiruan, kesegaran ikan, pengolahan
citra digital, ruang warna HSV
%-------------------------------------------------------------------
\end{spacing}

%==================================================================
% DAFTAR ISI PROPOSAL
% Acuan: Pedoman Tabel 6.3 - bagian awal proposal memuat Daftar Isi saja
% Daftar isi memakai jarak spasi 1
%==================================================================
\makeatletter

\clearpage
\chapter*{DAFTAR ISI}
\frontmatterentry{DAFTAR ISI}
\begin{spacing}{1}
\@starttoc{toc}
\end{spacing}

\makeatother


%==================================================================
% BAGIAN UTAMA PROPOSAL: BAB 1 - BAB 3
%==================================================================
\clearpage
\pagenumbering{arabic}
\setcounter{page}{1}

%==================================================================
% BAB 1 PENDAHULUAN
% Struktur mengikuti Pedoman 6.1 (Kerangka Proposal Tugas Akhir)
%  1.1 Latar Belakang
%  1.2 Perumusan Masalah
%  1.3 Tujuan
%  1.4 Manfaat Tugas Akhir
%  1.5 Batasan Masalah
%==================================================================
\chapter{PENDAHULUAN}
\label{chap:pendahuluan}

%------------------------------------------------------------------
\section{Latar Belakang}
%------------------------------------------------------------------

Indonesia merupakan negara penghasil ikan yang besar, sehingga mutu
produk perikanan menjadi faktor penting dalam rantai perdagangan.
Ikan termasuk bahan pangan yang sangat mudah mengalami kemunduran mutu
setelah ditangkap. Menurut \citet{estiasih2015}, proses kemunduran mutu
ikan berlangsung melalui tahap \textit{rigor mortis}, penguraian
senyawa oleh enzim, dan penguraian oleh aktivitas mikroorganisme.
Ketiga tahap tersebut berjalan cepat, terutama bila ikan tidak segera
didinginkan.

Penilaian kesegaran ikan di pasar tradisional umumnya dilakukan secara
organoleptik, yaitu dengan mengamati warna insang, mata, tekstur, dan
bau. Cara ini membutuhkan tenaga terlatih, bersifat subjektif, dan sulit
dilakukan dalam jumlah besar. Beberapa penelitian menunjukkan bahwa
penilaian berbasis citra digital dapat menjadi alternatif yang lebih
konsisten \citep{prasetyo2021}. Citra mata dan insang ikan menyimpan
informasi warna yang berkorelasi dengan tingkat kesegaran, sehingga
dapat diolah secara kuantitatif \citep{gonzalez2018}.

Pengolahan citra digital telah banyak dimanfaatkan untuk pemeriksaan
mutu bahan pangan secara tidak merusak. Namun, pemanfaatan teknologi
tersebut di lingkungan pasar ikan Banyuwangi belum banyak dilakukan,
padahal kebutuhan akan alat bantu penilaian yang cepat dan murah cukup
tinggi. Berdasarkan uraian tersebut, penelitian ini merancang bangun
sistem deteksi kesegaran ikan berbasis pengolahan citra digital dan
jaringan saraf tiruan yang diharapkan dapat membantu pembeli maupun
pedagang dalam menentukan mutu ikan secara cepat dan objektif.

%------------------------------------------------------------------
\section{Perumusan Masalah}
%------------------------------------------------------------------

Berdasarkan latar belakang di atas, rumusan masalah pada penelitian ini
adalah sebagai berikut.
\begin{packed_enum}
  \item Bagaimana merancang sistem akuisisi citra mata dan insang ikan
        yang menghasilkan gambar dengan pencahayaan konsisten?
  \item Bagaimana mengekstraksi ciri warna citra ikan pada ruang warna
        \textit{Hue, Saturation, Value} (HSV) sehingga dapat membedakan
        tingkat kesegaran ikan?
  \item Bagaimana kinerja jaringan saraf tiruan dalam mengklasifikasikan
        tingkat kesegaran ikan berdasarkan ciri warna yang diperoleh?
\end{packed_enum}

%------------------------------------------------------------------
\section{Tujuan}
%------------------------------------------------------------------

Tujuan yang ingin dicapai pada penelitian ini adalah:
\begin{packed_enum}
  \item merancang bangun sistem deteksi kesegaran ikan berbasis
        pengolahan citra digital;
  \item mengimplementasikan jaringan saraf tiruan untuk
        mengklasifikasikan tingkat kesegaran ikan; dan
  \item mengukur tingkat akurasi sistem yang dirancang pada data uji
        yang belum pernah digunakan pada proses pelatihan.
\end{packed_enum}

%------------------------------------------------------------------
\section{Manfaat Tugas Akhir}
%------------------------------------------------------------------

Manfaat yang diharapkan dari penelitian ini adalah sebagai berikut.
\begin{packed_enum}
  \item Bagi pembeli dan pedagang ikan, tersedia alat bantu penilaian
        kesegaran yang cepat, murah, dan tidak merusak produk.
  \item Bagi Program Studi \prodi, menjadi contoh penerapan pengolahan
        citra digital dan pembelajaran mesin pada permasalahan nyata.
  \item Bagi peneliti selanjutnya, menjadi rujukan awal pengembangan
        sistem sejenis dengan cakupan jenis ikan yang lebih luas.
\end{packed_enum}

%------------------------------------------------------------------
\section{Batasan Masalah}
%------------------------------------------------------------------

Agar pembahasan terarah, penelitian ini dibatasi pada hal-hal berikut.
\begin{packed_enum}
  \item Objek penelitian adalah ikan bandeng (\textit{Chanos chanos})
        yang diperoleh dari Pasar Banyuwangi.
  \item Ciri yang digunakan terbatas pada ciri warna pada ruang HSV,
        tidak mencakup ciri bentuk maupun tekstur.
  \item Pengambilan citra dilakukan di dalam kotak akusisi dengan
        pencahayaan buatan yang tetap.
  \item Klasifikasi dibagi menjadi tiga kelas, yaitu segar, cukup segar,
        dan tidak segar, mengacu pada kriteria organoleptik.
\end{packed_enum}

%==================================================================
% BAB 2 TINJAUAN PUSTAKA
% Berisi teori pendukung dan penelitian terdahulu.
% Perhatikan contoh penulisan rumus, tabel, dan sitasi di bab ini.
%==================================================================
\chapter{TINJAUAN PUSTAKA}
\label{chap:tinjauanpustaka}

%------------------------------------------------------------------
\section{Kesegaran Ikan}
%------------------------------------------------------------------

Kesegaran ikan adalah kondisi mutu ikan yang ditentukan oleh perubahan
biokimia dan mikrobiologis setelah ikan ditangkap. Menurut
\citet{estiasih2015}, kemunduran mutu ikan dipengaruhi oleh suhu
penyimpanan, cara penanganan, dan kandungan lemak pada tubuh ikan.
Semakin tinggi suhu penyimpanan, semakin cepat pula proses kemunduran
mutu berlangsung.

Secara organoleptik, kesegaran ikan dapat dinilai melalui penampakan
mata, warna insang, tekstur daging, dan bau. Ikan segar memiliki mata
yang jernih dan cembung, insang berwarna merah cerah, serta daging yang
kenyal. Seiring bertambahnya waktu penyimpanan, mata menjadi keruh dan
cekung, insang berubah pucat hingga kecoklatan, dan daging menjadi
lembek. Perubahan warna tersebut menjadi dasar pemilihan ciri warna
pada penelitian ini.

%------------------------------------------------------------------
\section{Pengolahan Citra Digital}
%------------------------------------------------------------------

Pengolahan citra digital merupakan proses mengolah citra menjadi citra
lain atau menjadi data yang dapat diinterpretasikan oleh komputer
\citep{gonzalez2018}. Citra digital direpresentasikan sebagai fungsi
dua dimensi $f(x,y)$ dengan $x$ dan $y$ menyatakan koordinat spasial
dan nilai fungsi menyatakan intensitas atau tingkat kecerahan piksel
pada koordinat tersebut.

Tahapan pengolahan citra yang digunakan pada penelitian ini meliputi
penyekalaan ukuran citra, perbaikan kontras, pemisahan kanal warna, dan
ekstraksi ciri statistik. Penyekalaan ukuran dilakukan agar seluruh
citra masukan memiliki dimensi yang seragam sebelum diproses lebih
lanjut.

%------------------------------------------------------------------
\section{Ruang Warna HSV}
%------------------------------------------------------------------

Ruang warna HSV memisahkan informasi warna (\textit{hue}) dari
informasi kecerahan (\textit{value}), sehingga lebih tahan terhadap
perubahan intensitas cahaya dibandingkan ruang warna RGB. Konversi
nilai RGB ke HSV dihitung menggunakan \cref{eq:hsv}.

\begin{equation}
\begin{aligned}
V &= \max(R,G,B) \\
S &= \begin{cases}
       \dfrac{V-\min(R,G,B)}{V}, & V \neq 0 \\[6pt]
       0, & V = 0
     \end{cases} \\
H &= \begin{cases}
       60^{\circ}\left(\dfrac{G-B}{V-\min(R,G,B)}\right), & V = R \\[6pt]
       60^{\circ}\left(2+\dfrac{B-R}{V-\min(R,G,B)}\right), & V = G \\[6pt]
       60^{\circ}\left(4+\dfrac{R-G}{V-\min(R,G,B)}\right), & V = B
     \end{cases}
\end{aligned}
\label{eq:hsv}
\end{equation}

\noindent dengan $R$, $G$, dan $B$ adalah nilai kanal merah, hijau, dan
biru yang telah dinormalisasi pada rentang 0 sampai 1. Nilai $H$
dinyatakan dalam derajat, sedangkan $S$ dan $V$ berada pada rentang 0
sampai 1.

%------------------------------------------------------------------
\section{Jaringan Saraf Tiruan}
%------------------------------------------------------------------

Jaringan saraf tiruan dengan algoritme perambatan balik
(\textit{backpropagation}) bekerja melalui dua tahap, yaitu perambatan
maju untuk menghitung keluaran dan perambatan balik untuk memperbaiki
bobot. Keluaran setiap neuron dihitung menggunakan fungsi aktivasi
sigmoid seperti pada \cref{eq:sigmoid}.

\begin{equation}
y_j = \frac{1}{1+e^{-\left(\sum_{i=1}^{n} w_{ij}x_i + b_j\right)}}
\label{eq:sigmoid}
\end{equation}

\noindent dengan $x_i$ adalah masukan ke-$i$, $w_{ij}$ adalah bobot dari
masukan ke-$i$ menuju neuron ke-$j$, dan $b_j$ adalah bias neuron ke-$j$.
Perbaikan bobot dilakukan dengan meminimumkan galat yang dinyatakan
pada \cref{eq:galat}.

\begin{equation}
E = \frac{1}{2}\sum_{j=1}^{m}\left(t_j-y_j\right)^2
\label{eq:galat}
\end{equation}

\noindent dengan $t_j$ adalah target keluaran dan $y_j$ adalah keluaran
jaringan. Bobot diperbarui sebesar $w_{ij} \leftarrow w_{ij} +
\eta\,\delta_j\,x_i$, dengan $\eta$ adalah laju pembelajaran dan
$\delta_j$ adalah gradien galat pada neuron ke-$j$.

%------------------------------------------------------------------
\section{Penelitian Terdahulu}
%------------------------------------------------------------------

Beberapa penelitian sebelumnya telah memanfaatkan citra digital untuk
menilai kesegaran ikan. Ringkasan penelitian tersebut disajikan pada
\cref{tab:terdahulu}.

\begin{table}[H]
\caption{Perbandingan Penelitian Terdahulu tentang Deteksi Kesegaran Ikan}
\label{tab:terdahulu}
\begin{spacing}{1}
\centering
\begin{tabularx}{\textwidth}{|>{\centering\arraybackslash}p{0.8cm}|X|>{\raggedright\arraybackslash}p{3.2cm}|>{\centering\arraybackslash}p{1.6cm}|}
\hline
\rowcolor{blue!12}
\textbf{No} & \textbf{Peneliti} & \textbf{Metode} & \textbf{Akurasi} \\
\hline
1 & \citet{prasetyo2021} & JST perambatan balik dengan ciri warna RGB & 88,00\% \\
\hline
2 & \citet{satria2004} & Analisis tekstur insang ikan & 84,50\% \\
\hline
3 & \citet{sudirman2010} & Klasifikasi warna mata ikan & 81,25\% \\
\hline
\end{tabularx}
\par\vspace{2pt}
\raggedright\noindent Sumber: hasil kajian pustaka (2024)
\end{spacing}
\end{table}

Berdasarkan \cref{tab:terdahulu}, sebagian besar penelitian menggunakan
ciri warna dan tekstur dengan akurasi di bawah 90\%. Penelitian ini
menggabungkan ciri warna pada ruang HSV dengan jaringan saraf tiruan
dua lapisan tersembunyi, yang diharapkan dapat meningkatkan akurasi
klasifikasi.

%==================================================================
% BAB 3 METODE PENELITIAN
% Berisi waktu dan tempat, alat dan bahan, tahapan penelitian,
% rancangan sistem, pengujian, dan jadwal pelaksanaan.
%==================================================================
\chapter{METODE PENELITIAN}
\label{chap:metode}

%------------------------------------------------------------------
\section{Waktu dan Tempat Pelaksanaan}
%------------------------------------------------------------------

Penelitian dilaksanakan pada bulan Februari sampai dengan Juni
\tahun. Pengambilan data citra dilakukan di Pasar Banyuwangi, sedangkan
proses pengolahan citra dan pelatihan jaringan saraf tiruan dilakukan
di Laboratorium Komputer Program Studi \prodi\ \institusi.

%------------------------------------------------------------------
\section{Alat dan Bahan}
%------------------------------------------------------------------

Alat dan bahan yang digunakan pada penelitian ini disajikan pada
\cref{tab:alatbahan}.

\begin{table}[H]
\caption{Alat dan Bahan Penelitian}
\label{tab:alatbahan}
\begin{spacing}{1}
\centering
\begin{tabularx}{\textwidth}{|>{\centering\arraybackslash}p{0.8cm}|X|>{\centering\arraybackslash}p{2.4cm}|}
\hline
\rowcolor{blue!12}
\textbf{No} & \textbf{Nama Alat/Bahan} & \textbf{Jumlah} \\
\hline
1 & Kamera digital 12 megapiksel & 1 unit \\
\hline
2 & Kotak akusisi citra berukuran 40 cm $\times$ 30 cm $\times$ 30 cm & 1 unit \\
\hline
3 & Lampu \textit{light emitting diode} (LED) 12 watt & 2 buah \\
\hline
4 & Ikan bandeng segar & 100 ekor \\
\hline
5 & Komputer dengan prosesor \textit{Intel Core i5} & 1 unit \\
\hline
6 & Perangkat lunak Python 3.11 dan OpenCV & 1 paket \\
\hline
\end{tabularx}
\par\vspace{2pt}
\raggedright\noindent Sumber: perencanaan penelitian (2024)
\end{spacing}
\end{table}

%------------------------------------------------------------------
\section{Tahapan Penelitian}
%------------------------------------------------------------------

Penelitian dilaksanakan melalui beberapa tahap yang saling berkaitan.
Tahapan tersebut digambarkan pada \cref{fig:tahapan}.

\begin{figure}[H]
\centering
\begin{tikzpicture}[node distance=0.5cm]
  \node (mulai)    [mulai]                        {Mulai};
  \node (studi)    [proses, below=of mulai]       {Studi literatur dan\\penyusunan instrumen};
  \node (akuisisi) [proses, below=of studi]       {Akuisisi citra mata\\dan insang ikan};
  \node (pra)      [proses, below=of akuisisi]    {Praproses: \textit{resize},\\perbaikan kontras};
  \node (ciri)     [proses, below=of pra]         {Ekstraksi ciri warna\\pada ruang HSV};
  \node (latih)    [proses, below=of ciri]        {Pelatihan jaringan\\saraf tiruan};
  \node (uji)      [proses, below=of latih]       {Pengujian model};
  \node (keputusan)[keputusan, below=of uji]      {Akurasi\\$\geq 90\%$?};
  \node (analisis) [proses, below=of keputusan]   {Analisis hasil dan\\penyusunan laporan};
  \node (selesai)  [mulai, below=of analisis]     {Selesai};

  \draw [panah] (mulai)    -- (studi);
  \draw [panah] (studi)    -- (akuisisi);
  \draw [panah] (akuisisi) -- (pra);
  \draw [panah] (pra)      -- (ciri);
  \draw [panah] (ciri)     -- (latih);
  \draw [panah] (latih)    -- (uji);
  \draw [panah] (uji)      -- (keputusan);
  \draw [panah] (keputusan) -- node[right, font=\small] {Ya} (analisis);
  \draw [panah] (analisis) -- (selesai);

  % Jalur kembali bila akurasi belum tercapai
  \draw [panah] (keputusan.east) -- ++(1.2,0)
        |- node[pos=0.25, right, font=\small] {Tidak} (latih.east);
\end{tikzpicture}
\caption{Diagram Alir Tahapan Penelitian}
\label{fig:tahapan}
\end{figure}

%------------------------------------------------------------------
\section{Perancangan Sistem}
%------------------------------------------------------------------

\subsection{Akuisisi Citra}

Pengambilan citra dilakukan di dalam kotak akusisi dengan dua lampu LED
yang dipasang pada sudut 45 derajat terhadap bidang objek. Jarak kamera
ke objek ditetapkan 25 cm dan seluruh citra diambil dengan pengaturan
ISO, kecepatan rana, serta panjang fokus yang sama agar intensitas
cahaya antar citra tetap seragam.

\subsection{Praproses Citra}

Citra hasil akuisisi disekala ulang menjadi $256 \times 256$ piksel
menggunakan interpolasi \textit{bilinear}. Perbaikan kontras dilakukan
dengan perataan histogram untuk memperjelas perbedaan warna pada mata
dan insang ikan.

\subsection{Ekstraksi Ciri Warna}

Ekstraksi ciri dilakukan pada ruang warna HSV menggunakan
\cref{eq:hsv}. Nilai rata-rata dan simpangan baku setiap kanal
dihitung dan dijadikan satu vektor ciri berukuran enam dimensi, yaitu
$\bar{H}$, $\bar{S}$, $\bar{V}$, $s_H$, $s_S$, dan $s_V$.

\subsection{Klasifikasi}

Vektor ciri dinormalisasi ke rentang 0 sampai 1, kemudian digunakan
sebagai masukan jaringan saraf tiruan dengan arsitektur 6 neuron
masukan, dua lapisan tersembunyi berisi 16 dan 8 neuron, serta 3 neuron
keluaran yang mewakili kelas segar, cukup segar, dan tidak segar.

%------------------------------------------------------------------
\section{Pengujian dan Evaluasi}
%------------------------------------------------------------------

Pengujian dilakukan menggunakan \textit{k-fold cross validation} dengan
$k = 5$. Kinerja sistem diukur menggunakan akurasi seperti pada
\cref{eq:akurasi}.

\begin{equation}
\mathrm{Akurasi} = \frac{\text{jumlah data uji yang benar}}{\text{jumlah seluruh data uji}} \times 100\%
\label{eq:akurasi}
\end{equation}

\noindent Selain akurasi, dihitung pula presisi dan perolehan
(\textit{recall}) untuk setiap kelas menggunakan
\cref{eq:presisi,eq:recall}.

\begin{equation}
\mathrm{Presisi} = \frac{TP}{TP+FP} \times 100\%
\label{eq:presisi}
\end{equation}

\begin{equation}
\mathrm{Perolehan} = \frac{TP}{TP+FN} \times 100\%
\label{eq:recall}
\end{equation}

\noindent dengan $TP$ adalah \textit{true positive}, $FP$ adalah
\textit{false positive}, dan $FN$ adalah \textit{false negative}.

%------------------------------------------------------------------
\section{Jadwal Pelaksanaan}
%------------------------------------------------------------------

Jadwal pelaksanaan penelitian disusun dalam bentuk \textit{bar-chart}
seperti pada \cref{tab:jadwal}.

\begin{table}[H]
\caption{Model \textit{Bar-Chart} Penyusunan Jadwal Pelaksanaan Tugas Akhir}
\label{tab:jadwal}
\begin{spacing}{1}
\centering
\begin{tabular}{|>{\raggedright\arraybackslash}p{5.2cm}|*{6}{>{\centering\arraybackslash}p{0.9cm}|}}
\hline
\multirow{2}{*}{\textbf{Jenis Kegiatan}} & \multicolumn{6}{c|}{\textbf{Bulan Ke-}} \\
\cline{2-7}
 & \textbf{1} & \textbf{2} & \textbf{3} & \textbf{4} & \textbf{5} & \textbf{6} \\
\hline
Studi literatur            & \cellcolor{blue!25} & & & & & \\
\hline
Penyusunan instrumen       & \cellcolor{blue!25} & \cellcolor{blue!25} & & & & \\
\hline
Akuisisi citra             & & \cellcolor{blue!25} & & & & \\
\hline
Praproses dan ekstraksi    & & & \cellcolor{blue!25} & & & \\
\hline
Pelatihan dan pengujian    & & & & \cellcolor{blue!25} & & \\
\hline
Analisis dan penulisan     & & & & & \cellcolor{blue!25} & \cellcolor{blue!25} \\
\hline
\end{tabular}
\par\vspace{2pt}
\raggedright\noindent Sumber: perencanaan penelitian (2024)
\end{spacing}
\end{table}


%==================================================================
% BAGIAN AKHIR
%==================================================================
%==================================================================
% DAFTAR PUSTAKA
% Gaya sitasi: APA (penulis, tahun) - Pedoman 6.1 Bagian Akhir
%==================================================================
\clearpage
\renewcommand{\bibname}{DAFTAR PUSTAKA}
\frontmatterentry{DAFTAR PUSTAKA}

\begin{spacing}{1}
\bibliography{a9-pustaka}
\end{spacing}

%==================================================================
% LAMPIRAN
% Setiap lampiran dibuat dengan perintah \lampiran{Judul Lampiran}
% dan otomatis muncul di DAFTAR LAMPIRAN.
%==================================================================

\lampiran{Kode Program Pelatihan Jaringan Saraf Tiruan}

\noindent Berikut adalah kode program lengkap untuk melatih model
jaringan saraf tiruan yang digunakan pada penelitian ini.

\lstinputlisting[language=Python,
                 caption={Kode Program Pelatihan Model},
                 label={lst:pelatihan}]{kode/latih_model.py}

\clearpage

\lampiran{Lembar Observasi Kesegaran Ikan}

\noindent Formulir berikut digunakan untuk mencatat hasil pengamatan
organoleptik pada setiap objek penelitian.

\begin{spacing}{1}
\begin{tabularx}{\textwidth}{|>{\centering\arraybackslash}p{0.9cm}|X|>{\centering\arraybackslash}p{2.8cm}|>{\centering\arraybackslash}p{2.8cm}|}
\hline
\rowcolor{blue!12}
\textbf{No} & \textbf{Parameter Pengamatan} & \textbf{Skor (1--3)} & \textbf{Keterangan} \\
\hline
1 & Kejernihan mata ikan & & \\
\hline
2 & Warna insang & & \\
\hline
3 & Kekenyalan daging & & \\
\hline
4 & Bau & & \\
\hline
5 & Permukaan tubuh & & \\
\hline
\end{tabularx}
\end{spacing}

\vspace{2pt}
\noindent Keterangan skor: 1 = tidak segar, 2 = cukup segar, 3 = segar.


\end{document}
